%% kapa: exam=1 round=23 year=2012 period=2 form=A subject=accounting subject_ko=회계학 seq=37 number=37
%% source: Q-Net 감정평가사 기출문제 (cst00309) artlSeq=5052580 file=raw/1008/5052580/2012년 제23회 감정평가사 제1차시험 기출문제(2교시 A형).pdf
%% conversion: kordoc 4.24.51 -> tools/gen1.py (auto); review=proofread-sonnet; flags=none
\begin{question}{37}
화장품을 제조하는 (주)한국의 20×1년 직접재료예산과 관련된 자료를 이용하여 계산한 3분기의 재료구입 예산액은?
\begin{framedtext}
ㆍ매출계획 및 재고계획에 따른 각 분기별 화장품 추정생산량 \par
\begin{center}\small
\begin{tabular}{|c|c|c|c|}
\hline
\multicolumn{4}{|c|}{20×1년} \\
\hline
1분기 & 2분기 & 3분기 & 4분기 \\
\hline
1,400단위 & 3,200단위 & 3,600단위 & 1,900단위 \\
\hline
\end{tabular}
\end{center}
ㆍ화장품 1단위를 생산하는데 필요한 재료는 15kg이다. \par
ㆍ분기말 목표재료재고량은 다음 분기 생산량에 필요한 재료의 10\%로 한다. \par
ㆍ재료구입단가는 kg당 ￦2이다. \par
\end{framedtext}
\begin{choices}
  \item ￦102,900
  \item ￦103,200
  \item ￦104,100
  \item ￦105,300
  \item ￦106,400
\end{choices}
\end{question}
\answer{1}
